% fig-eval-attacktree.tex — attack routes and their detection controls.
% Left group: A1/A3/A5 are independent substitution routes to the goal (an OR:
% any one suffices); each attack box sits above the control that detects it.
% Right group: A2/A4 are detection-evasion attempts, each also detected.
% Dashed arrows read "detected by". X1 (endpoint compromise) bypasses the
% directory entirely and is out of scope, shown separately below.
% Label: \label{fig:eval-attacktree}
\resizebox{\textwidth}{!}{%
\begin{tikzpicture}[>=Stealth,font=\scriptsize,
  node distance=0pt,
  goal/.style={draw,thick,rounded corners,fill=black!6,align=center,text width=3.6cm,minimum height=9mm,inner sep=3pt},
  gate/.style={draw,rounded corners,fill=blue!8,draw=blue!55!black,align=center,font=\scriptsize\bfseries,text width=3.6cm,minimum height=7mm,inner sep=3pt},
  atk/.style={draw,rounded corners,draw=red!60!black,align=center,text width=2.4cm,minimum height=10mm,inner sep=3pt},
  det/.style={draw,rounded corners,draw=green!45!black,fill=green!5,text=green!25!black,align=center,text width=2.4cm,minimum height=12mm,inner sep=3pt},
  grp/.style={font=\scriptsize\itshape,text=black!60,align=center},
  oos/.style={draw,dashed,rounded corners,draw=orange!80!black,align=center,text width=12.2cm,minimum height=9mm,inner sep=4pt}]

  % ---- left group: substitution routes (A1, A3, A5) ----
  \node[goal] (g)   at (-2.0, 5.3) {Goal: client accepts an attacker-controlled key};
  \node[gate] (orG) at (-2.0, 3.7) {OR: any one substitution route suffices};
  \draw[->,black!45] (g) -- (orG);

  \node[atk] (a1) at (-4.7, 2.0) {A1\\server substitution};
  \node[atk] (a3) at (-2.0, 2.0) {A3\\network MitM};
  \node[atk] (a5) at ( 0.7, 2.0) {A5\\malicious rotation};
  \foreach \n in {a1,a3,a5}{\draw[->,black!45] (orG) -- (\n);}

  \node[det] (d1) at (-4.7, -0.2) {pinned-key self-check};
  \node[det] (d3) at (-2.0, -0.2) {honest-directory VRF comparison};
  \node[det] (d5) at ( 0.7, -0.2) {re-fetch and predecessor cross-signature};

  % ---- right group: detection-evasion attempts (A2, A4) ----
  \node[grp] (rhead) at (4.7, 3.7) {Detection-evasion attempts};
  \node[atk] (a2) at (3.4, 2.0) {A2\\warning suppression};
  \node[atk] (a4) at (6.0, 2.0) {A4\\equivocation (split view)};
  \draw[->,black!45] (rhead) -- (a2);
  \draw[->,black!45] (rhead) -- (a4);

  \node[det] (d2) at (3.4, -0.2) {mismatch and proven warning absence};
  \node[det] (d4) at (6.0, -0.2) {cross-vantage fork check};

  % ---- detected-by arrows (uniform, dashed) ----
  \foreach \a/\d in {a1/d1,a3/d3,a5/d5,a2/d2,a4/d4}{%
    \draw[->,green!45!black,dashed] (\a) -- (\d);}

  \node[grp,anchor=west] at (-5.9, -1.4) {Dashed arrow: the control that detects the attack.};

  % ---- X1: out of scope, shown separately ----
  \node[oos] (x1) at (0.05, -2.4) {X1 (endpoint compromise): the attacker uses the victim's own
    device and key, bypassing the directory entirely, so it reaches no directory control and is not
    evaluated here (MITRE~T1552.004).};
\end{tikzpicture}%
}
